\documentclass[11pt, a4paper]{ctexart}

\usepackage{assignment_style}

\begin{document}

% =========================================
% 1
% =========================================
\begin{problem}{1. 设插值节点为 $x_i=-1,0,1$，函数值为 $f(x_i)=2^{x_i}$。}
    \begin{enumerate}[(1)]
        \item 用线性插值计算 $2^{0.3}$ 的近似值，并估计误差；
        \item 用二次插值计算 $2^{0.3}$ 的近似值，并估计误差。
    \end{enumerate}
\end{problem}

\begin{solution}
    由
    \[
        f(-1)=\frac{1}{2},\qquad f(0)=1,\qquad f(1)=2.
    \]

    \textbf{(1) 线性插值}

    在区间 $[0,1]$ 上作一次 Lagrange 插值，
    \[
        L_1(x)=f(0)\frac{x-1}{0-1}+f(1)\frac{x-0}{1-0}=x+1.
    \]
    因而
    \[
        L_1(0.3)=1.3.
    \]
    余项为
    \[
        R_1(x)=\frac{f''(\xi)}{2!}(x-0)(x-1),\qquad \xi\in(0,1).
    \]
    由于
    \[
        f''(x)=2^x(\ln 2)^2,
    \]
    取 $\xi=0.5$ 估计，则
    \[
        |R_1(0.3)|\approx \left|\frac{2^{0.5}(\ln 2)^2}{2}\cdot 0.3\cdot(-0.7)\right|
        =0.071344.
    \]

    \textbf{(2) 二次插值}

    取节点 $-1,0,1$，二次 Lagrange 插值多项式为
    \[
        L_2(x)=\frac{1}{2}\cdot\frac{x(x-1)}{(-1-0)(-1-1)}
        +1\cdot\frac{(x+1)(x-1)}{(0+1)(0-1)}
        +2\cdot\frac{x(x+1)}{(1+1)(1-0)}.
    \]
    化简得
    \[
        L_2(x)=1+\frac{3}{4}x+\frac{1}{4}x^2.
    \]
    因而
    \[
        L_2(0.3)=1+\frac{3}{4}\cdot 0.3+\frac{1}{4}\cdot 0.3^2=1.2475.
    \]
    余项为
    \[
        R_2(x)=\frac{f^{(3)}(\xi)}{3!}(x+1)x(x-1),\qquad \xi\in(-1,1),
    \]
    而
    \[
        f^{(3)}(x)=2^x(\ln 2)^3.
    \]
    取 $\xi=0$ 估计，则
    \[
        |R_2(0.3)|\approx \left|\frac{2^0(\ln 2)^3}{6}\cdot 1.3\cdot 0.3\cdot(-0.7)\right|
        =0.015153.
    \]
\end{solution}


% =========================================
% 3
% =========================================
\begin{problem}{3. 已知插值节点}
    \[
        x_i: 1,\ 3,\ 4,\ 6,\qquad f(x_i): -7,\ 5,\ 8,\ 14,
    \]
    求三次 Lagrange 插值多项式，并计算 $f(2)$ 的近似值。
\end{problem}

\begin{solution}
    三次 Lagrange 插值多项式为
    \[
        L_3(x)=\sum_{i=0}^{3} f(x_i)\ell_i(x),
    \]
    计算后可化为
    \[
        L_3(x)=\frac{1}{5}x^3-\frac{13}{5}x^2+\frac{69}{5}x-\frac{92}{5}.
    \]
    代入 $x=2$ 得
    \[
        L_3(2)=\frac{8}{5}-\frac{52}{5}+\frac{138}{5}-\frac{92}{5}=\frac{2}{5}=0.4.
    \]
    因而
    \[
        f(2)\approx 0.4.
    \]
\end{solution}


% =========================================
% 6
% =========================================
\begin{problem}{6. 已知节点}
    \[
        x_i: 1,\ 1.1,\ 1.2,\ 1.3,\ 1.4,\qquad f(x_i)=e^{x_i^2-1}.
    \]
    \begin{enumerate}[(1)]
        \item 计算四次插值多项式在 $x=1.25$ 处的值，并估计误差；
        \item 求 $f(x)$ 的四次 Newton 插值多项式，并计算 $f(1.25)$；
        \item 增加节点 $x_5=1.5$ 后，求五次 Newton 插值多项式在 $x=1.25$ 处的值。
    \end{enumerate}
\end{problem}

\begin{solution}
    记
    \[
        f(x)=e^{x^2-1}.
    \]

    \textbf{(1) 四次插值值与误差估计}

    由插值计算可得
    \[
        L_4(1.25)\approx 1.7550.
    \]
    四次插值余项为
    \[
        R_4(x)=\frac{f^{(5)}(\xi)}{5!}\prod_{i=0}^{4}(x-x_i),\qquad \xi\in[1,1.4].
    \]
    对本题，
    \[
        f^{(5)}(x)=e^{x^2-1}(120x+160x^3+32x^5).
    \]
    当 $x=1.25$ 时，
    \[
        \frac{\left|\prod_{i=0}^{4}(1.25-x_i)\right|}{5!}
        =1.171875\times 10^{-7}.
    \]
    取 $\xi=1.2$ 估计，可得
    \[
        |R_4(1.25)|\approx 9.0998\times 10^{-5}.
    \]

    \textbf{(2) 四次 Newton 插值}

    由差商表可得首行系数
    \[
        1.0000,\ 2.3368,\ 4.2676,\ 6.1047,\ 7.6523.
    \]
    因而四次 Newton 插值多项式可写为
    \[
        \begin{aligned}
            P_4(x)=&\,1
            +2.3368(x-1)
            +4.2676(x-1)(x-1.1) \\
            &+6.1047(x-1)(x-1.1)(x-1.2) \\
            &+7.6523(x-1)(x-1.1)(x-1.2)(x-1.3).
        \end{aligned}
    \]
    代入 $x=1.25$，
    \[
        P_4(1.25)\approx 1.7550.
    \]

    \textbf{(3) 增加节点 $x_5=1.5$ 后}

    增加节点后，差商表首行新增系数 $8.6117$，因此
    \[
        P_5(x)=P_4(x)+8.6117(x-1)(x-1.1)(x-1.2)(x-1.3)(x-1.4).
    \]
    从而
    \[
        P_5(1.25)\approx 1.7551.
    \]
\end{solution}


% =========================================
% 7
% =========================================
\begin{problem}{7. 设 $x_0,\ldots,x_n$ 为 $n+1$ 个互异节点，$\ell_j(x)\ (j=0,\ldots,n)$ 为这组节点上的 Lagrange 插值基函数。证明：}
    \begin{enumerate}[(1)]
        \item 对任意 $k=1,\ldots,n$，有
        \[
            \sum_{j=0}^{n}(x_j-x)^k\ell_j(x)=0;
        \]
        \item 若 $p(x)$ 是首项系数为 $1$ 的 $n+1$ 次多项式，则
        \[
            p(x)-\sum_{j=0}^{n}p(x_j)\ell_j(x)=w_{n+1}(x):=\prod_{j=0}^{n}(x-x_j).
        \]
    \end{enumerate}
\end{problem}

\begin{myproof}
    \textbf{(1)} 令
    \[
        f(t)=(t-x)^k,\qquad 1\le k\le n.
    \]
    这是一个次数不超过 $n$ 的多项式，因此其在节点 $x_0,\ldots,x_n$ 上的 Lagrange 插值多项式就是其本身，即
    \[
        f(t)=\sum_{j=0}^{n}f(x_j)\ell_j(t)
        =\sum_{j=0}^{n}(x_j-x)^k\ell_j(t).
    \]
    令 $t=x$，则左端为 $(x-x)^k=0$，故
    \[
        \sum_{j=0}^{n}(x_j-x)^k\ell_j(x)=0.
    \]

    \textbf{(2)} 令
    \[
        L_n(x)=\sum_{j=0}^{n}p(x_j)\ell_j(x)
    \]
    为 $p(x)$ 在节点 $x_0,\ldots,x_n$ 上的 Lagrange 插值多项式。因为 $p(x)$ 是 $n+1$ 次多项式，故插值余项满足
    \[
        p(x)-L_n(x)=\frac{p^{(n+1)}(\xi)}{(n+1)!}\prod_{j=0}^{n}(x-x_j).
    \]
    又因为 $p(x)$ 的首项系数为 $1$，所以
    \[
        p^{(n+1)}(\xi)=(n+1)!.
    \]
    因而
    \[
        p(x)-\sum_{j=0}^{n}p(x_j)\ell_j(x)
        =\prod_{j=0}^{n}(x-x_j)
        =w_{n+1}(x).
    \]
    命题得证。
\end{myproof}


% =========================================
% 8
% =========================================
\begin{problem}{8. 已知节点}
    \[
        x_i: 1.615,\ 1.634,\ 1.702,\ 1.828,
    \]
    \[
        f(x_i): 2.41450,\ 2.46259,\ 2.65271,\ 3.03035.
    \]
    求三次 Newton 插值多项式，并计算 $f(1.682)$ 的近似值。
\end{problem}

\begin{solution}
    差商表首行系数为
    \[
        2.41450,\qquad 2.5311,\qquad 3.0440,\qquad -9.4206.
    \]
    因而三次 Newton 插值多项式为
    \[
        \begin{aligned}
            P_3(x)=&\,2.41450+2.5311(x-1.615) \\
            &+3.0440(x-1.615)(x-1.634) \\
            &-9.4206(x-1.615)(x-1.634)(x-1.702).
        \end{aligned}
    \]
    代入 $x=1.682$，得
    \[
        P_3(1.682)\approx 2.5945.
    \]
    因而
    \[
        f(1.682)\approx 2.5945.
    \]
\end{solution}


% =========================================
% 15
% =========================================
\begin{problem}{15. 设 $f(x)=3xe^x-e^{2x}$，在节点 $x_0=1,\ x_1=1.05$ 上作三次 Hermite 插值，求插值多项式，并计算 $f(1.03)$ 的近似值及误差估计。}
\end{problem}

\begin{solution}
    先计算函数值与导数值：
    \[
        f(1)=0.76579,\qquad f(1.05)=0.83543,
    \]
    \[
        f'(x)=3e^x+3xe^x-2e^{2x},
    \]
    \[
        f'(1)=1.5316,\qquad f'(1.05)=1.2422.
    \]
    用重节点差商构造三次 Hermite 插值。由
    \[
        f[1,1]=1.5316,\qquad f[1,1,1.05]=-2.775,
    \]
    \[
        f[1,1,1.05,1.05]=-4.7502,
    \]
    可得
    \[
        H_3(x)=0.76579+1.5316(x-1)-2.775(x-1)^2-4.7502(x-1)^2(x-1.05).
    \]
    于是
    \[
        H_3(1.03)\approx 0.80932.
    \]

    Hermite 余项为
    \[
        R_3(x)=\frac{f^{(4)}(\xi)}{4!}(x-1)^2(x-1.05)^2,\qquad \xi\in(1,1.05),
    \]
    而
    \[
        f^{(4)}(x)=12e^x-16e^{2x}+3xe^x.
    \]
    取 $\xi=1.025$ 估计，则
    \[
        |R_3(1.03)|\approx 1.2341\times 10^{-6}.
    \]
\end{solution}


% =========================================
% 16
% =========================================
\begin{problem}{16. 设 $f(x)\in C^5[0,2]$，且满足}
    \[
        f(0)=1,\qquad f(1)=2,\qquad f(2)=1,
    \]
    \[
        f'(1)=0,\qquad f'(2)=-1.
    \]
    求四次 Hermite 插值多项式 $H_4(x)$，并写出误差项 $f(x)-H_4(x)$ 的表达式。
\end{problem}

\begin{solution}
    采用 Newton-Hermite 形式，重节点取为
    \[
        0,\ 1,\ 1,\ 2,\ 2.
    \]
    由差商表可得相应系数
    \[
        1,\qquad 1,\qquad -1,\qquad 0,\qquad \frac12.
    \]
    因而
    \[
        \begin{aligned}
            H_4(x)=&\,1+(x-0)-(x-0)(x-1) \\
            &+0\cdot (x-0)(x-1)^2
            +\frac12(x-0)(x-1)^2(x-2).
        \end{aligned}
    \]
    展开后得
    \[
        H_4(x)=1+x+\frac32x^2-2x^3+\frac12x^4.
    \]

    由于节点 $x=1,2$ 各重复一次，故误差项为
    \[
        f(x)-H_4(x)=\frac{f^{(5)}(\xi)}{5!}\,x(x-1)^2(x-2)^2,
        \qquad \xi\in(0,2).
    \]
\end{solution}

\end{document}
